% Options for packages loaded elsewhere
\PassOptionsToPackage{unicode}{hyperref}
\PassOptionsToPackage{hyphens}{url}
\PassOptionsToPackage{dvipsnames,svgnames,x11names}{xcolor}
\documentclass[
  11pt,
]{article}
\usepackage{xcolor}
\usepackage[margin=1in]{geometry}
\usepackage{amsmath,amssymb}
\setcounter{secnumdepth}{-\maxdimen} % remove section numbering
\usepackage{iftex}
\ifPDFTeX
  \usepackage[T1]{fontenc}
  \usepackage[utf8]{inputenc}
  \usepackage{textcomp} % provide euro and other symbols
\else % if luatex or xetex
  \usepackage{unicode-math} % this also loads fontspec
  \defaultfontfeatures{Scale=MatchLowercase}
  \defaultfontfeatures[\rmfamily]{Ligatures=TeX,Scale=1}
\fi
\usepackage{lmodern}
\ifPDFTeX\else
  % xetex/luatex font selection
\fi
% Use upquote if available, for straight quotes in verbatim environments
\IfFileExists{upquote.sty}{\usepackage{upquote}}{}
\IfFileExists{microtype.sty}{% use microtype if available
  \usepackage[]{microtype}
  \UseMicrotypeSet[protrusion]{basicmath} % disable protrusion for tt fonts
}{}
\makeatletter
\@ifundefined{KOMAClassName}{% if non-KOMA class
  \IfFileExists{parskip.sty}{%
    \usepackage{parskip}
  }{% else
    \setlength{\parindent}{0pt}
    \setlength{\parskip}{6pt plus 2pt minus 1pt}}
}{% if KOMA class
  \KOMAoptions{parskip=half}}
\makeatother
\usepackage{color}
\usepackage{fancyvrb}
\newcommand{\VerbBar}{|}
\newcommand{\VERB}{\Verb[commandchars=\\\{\}]}
\DefineVerbatimEnvironment{Highlighting}{Verbatim}{commandchars=\\\{\}}
% Add ',fontsize=\small' for more characters per line
\newenvironment{Shaded}{}{}
\newcommand{\AlertTok}[1]{\textcolor[rgb]{1.00,0.00,0.00}{\textbf{#1}}}
\newcommand{\AnnotationTok}[1]{\textcolor[rgb]{0.38,0.63,0.69}{\textbf{\textit{#1}}}}
\newcommand{\AttributeTok}[1]{\textcolor[rgb]{0.49,0.56,0.16}{#1}}
\newcommand{\BaseNTok}[1]{\textcolor[rgb]{0.25,0.63,0.44}{#1}}
\newcommand{\BuiltInTok}[1]{\textcolor[rgb]{0.00,0.50,0.00}{#1}}
\newcommand{\CharTok}[1]{\textcolor[rgb]{0.25,0.44,0.63}{#1}}
\newcommand{\CommentTok}[1]{\textcolor[rgb]{0.38,0.63,0.69}{\textit{#1}}}
\newcommand{\CommentVarTok}[1]{\textcolor[rgb]{0.38,0.63,0.69}{\textbf{\textit{#1}}}}
\newcommand{\ConstantTok}[1]{\textcolor[rgb]{0.53,0.00,0.00}{#1}}
\newcommand{\ControlFlowTok}[1]{\textcolor[rgb]{0.00,0.44,0.13}{\textbf{#1}}}
\newcommand{\DataTypeTok}[1]{\textcolor[rgb]{0.56,0.13,0.00}{#1}}
\newcommand{\DecValTok}[1]{\textcolor[rgb]{0.25,0.63,0.44}{#1}}
\newcommand{\DocumentationTok}[1]{\textcolor[rgb]{0.73,0.13,0.13}{\textit{#1}}}
\newcommand{\ErrorTok}[1]{\textcolor[rgb]{1.00,0.00,0.00}{\textbf{#1}}}
\newcommand{\ExtensionTok}[1]{#1}
\newcommand{\FloatTok}[1]{\textcolor[rgb]{0.25,0.63,0.44}{#1}}
\newcommand{\FunctionTok}[1]{\textcolor[rgb]{0.02,0.16,0.49}{#1}}
\newcommand{\ImportTok}[1]{\textcolor[rgb]{0.00,0.50,0.00}{\textbf{#1}}}
\newcommand{\InformationTok}[1]{\textcolor[rgb]{0.38,0.63,0.69}{\textbf{\textit{#1}}}}
\newcommand{\KeywordTok}[1]{\textcolor[rgb]{0.00,0.44,0.13}{\textbf{#1}}}
\newcommand{\NormalTok}[1]{#1}
\newcommand{\OperatorTok}[1]{\textcolor[rgb]{0.40,0.40,0.40}{#1}}
\newcommand{\OtherTok}[1]{\textcolor[rgb]{0.00,0.44,0.13}{#1}}
\newcommand{\PreprocessorTok}[1]{\textcolor[rgb]{0.74,0.48,0.00}{#1}}
\newcommand{\RegionMarkerTok}[1]{#1}
\newcommand{\SpecialCharTok}[1]{\textcolor[rgb]{0.25,0.44,0.63}{#1}}
\newcommand{\SpecialStringTok}[1]{\textcolor[rgb]{0.73,0.40,0.53}{#1}}
\newcommand{\StringTok}[1]{\textcolor[rgb]{0.25,0.44,0.63}{#1}}
\newcommand{\VariableTok}[1]{\textcolor[rgb]{0.10,0.09,0.49}{#1}}
\newcommand{\VerbatimStringTok}[1]{\textcolor[rgb]{0.25,0.44,0.63}{#1}}
\newcommand{\WarningTok}[1]{\textcolor[rgb]{0.38,0.63,0.69}{\textbf{\textit{#1}}}}
\usepackage{longtable,booktabs,array}
\newcounter{none} % for unnumbered tables
\usepackage{calc} % for calculating minipage widths
% Correct order of tables after \paragraph or \subparagraph
\usepackage{etoolbox}
\makeatletter
\patchcmd\longtable{\par}{\if@noskipsec\mbox{}\fi\par}{}{}
\makeatother
% Allow footnotes in longtable head/foot
\IfFileExists{footnotehyper.sty}{\usepackage{footnotehyper}}{\usepackage{footnote}}
\makesavenoteenv{longtable}
\setlength{\emergencystretch}{3em} % prevent overfull lines
\providecommand{\tightlist}{%
  \setlength{\itemsep}{0pt}\setlength{\parskip}{0pt}}
\usepackage{bookmark}
\IfFileExists{xurl.sty}{\usepackage{xurl}}{} % add URL line breaks if available
\urlstyle{same}
\hypersetup{
  pdftitle={UQFF 2027-2028 Quadruple Falsifier: Single-Parameter Joint Lock of P6{,} P10{,} P11{,} P12},
  pdfauthor={Daniel Murphy},
  colorlinks=true,
  linkcolor={Maroon},
  filecolor={Maroon},
  citecolor={Blue},
  urlcolor={Blue},
  pdfcreator={LaTeX via pandoc}}

\title{UQFF 2027-2028 Quadruple Falsifier: Single-Parameter Joint Lock
of P6, P10, P11, P12}
\author{Daniel Murphy}
\date{May 2026}

\begin{document}
\maketitle

\section{Abstract}\label{abstract}

PAPER\_1177 (Session 259) showed that three independent 2027 decisive
experiments --- Wuhan torsion-balance Kaluza-Klein tower search (P6),
LIGO O5 black-hole ringdown spectral offset (P11), and Euclid Y1
weak-lensing \(\sigma_8\) (P12) --- are unified by a single
dimensionless ratio \(\xi \equiv D_{\mathrm{crit}}/D_{\mathrm{BSFG}}\)
with UQFF canonical value \(\xi_0 = 13/3 \approx 4.333\). This paper
extends the joint chi-squared self-lock to a fourth independent
experiment: IceCube-Gen2 neutrino Cherenkov spectral cutoff (P10,
PAPER\_1163). All four are 2027-2028 measurements with mutually
orthogonal systematics. UQFF is required, under closed Lagrangian
self-consistency, to satisfy all four within \(1\sigma\) at the same
\(\xi\). Any single \(> 3\sigma\) outlier falsifies the entire R26
closure framework.

\section{1. Quadruple Joint chi\^{}2(xi)}\label{quadruple-joint-chi2xi}

The 2027-2028 falsifier set is

\[
\chi^2(\xi) = \sum_{k \in \{P6, P10, P11, P12\}} \frac{[O_k - M_k(\xi)]^2}{\sigma_k^2}
\]

with predictions

{\def\LTcaptype{none} % do not increment counter
\begin{longtable}[]{@{}
  >{\raggedleft\arraybackslash}p{(\linewidth - 8\tabcolsep) * \real{0.0818}}
  >{\raggedright\arraybackslash}p{(\linewidth - 8\tabcolsep) * \real{0.2000}}
  >{\raggedright\arraybackslash}p{(\linewidth - 8\tabcolsep) * \real{0.3455}}
  >{\raggedright\arraybackslash}p{(\linewidth - 8\tabcolsep) * \real{0.2727}}
  >{\raggedright\arraybackslash}p{(\linewidth - 8\tabcolsep) * \real{0.1000}}@{}}
\toprule\noalign{}
\begin{minipage}[b]{\linewidth}\raggedleft
Channel
\end{minipage} & \begin{minipage}[b]{\linewidth}\raggedright
Observable
\end{minipage} & \begin{minipage}[b]{\linewidth}\raggedright
Model \(M_k(\xi)\)
\end{minipage} & \begin{minipage}[b]{\linewidth}\raggedright
UQFF central \((\xi = \xi_0)\)
\end{minipage} & \begin{minipage}[b]{\linewidth}\raggedright
\(1\sigma\)
\end{minipage} \\
\midrule\noalign{}
\endhead
\bottomrule\noalign{}
\endlastfoot
P6 & \(L_{KK}\) (\(\mu\)m) & \(26.3 \cdot (\xi_0/\xi)\) & \(26.3\) &
\(0.5\) \\
P10 & \(E_{cut}\) (PeV) & \(6.3 \cdot (\xi/\xi_0)^{1/2}\) & \(6.3\) &
\(0.4\) \\
P11 & \(R_{21/22}\) & \(0.10 \cdot \xi^{0.25}\) & \(0.144\) &
\(0.010\) \\
P12 & \(\sigma_8\) & \(0.7851 \cdot (1 + 0.01525 \cdot \xi/\xi_0)\) &
\(0.797\) & \(0.005\) \\
\end{longtable}
}

\section{2. Joint chi\^{}2 threshold}\label{joint-chi2-threshold}

For 4 channels and 1 free parameter \(\xi\), degrees of freedom
\(\nu = 3\). The 95\% confidence threshold is \(\chi^2_{95}(3) = 7.81\).
The 99.7\% (3-sigma) threshold is \(\chi^2_{99.7}(3) = 14.16\).

UQFF is \textbf{decisively falsified} when \(\chi^2_{\min} > 14.16\) at
the best-fit \(\xi\).

\section{3. Cross-Lock Geometry}\label{cross-lock-geometry}

The four channels probe four different scales:

\begin{itemize}
\tightlist
\item
  P6: short-distance (\(\mu\)m) gravitational deformation.
\item
  P10: high-energy (\(\sim\)PeV) neutrino dispersion.
\item
  P11: stellar-mass (\(\sim 60 M_\odot\)) ringdown quasi-normal mode
  ratio.
\item
  P12: cosmological (\(\sim 100\) Mpc) matter power.
\end{itemize}

Each scale couples to a different aspect of the R26 manifold: KK radius,
vortex spectral index, BSFG bridging coefficient, and large-scale
clustering amplitude. UQFF's closed Lagrangian fixes a single \(\xi\)
across all four. Any disagreement is a closed-form refutation; no
nuisance parameter can absorb it.

\section{4. Decision Tree}\label{decision-tree}

\begin{verbatim}
                  [2027-2028 Quadruple Window]
                            |
        chi^2_min(xi_best_fit) computed from {P6, P10, P11, P12}
                            |
                +-----------+-----------+
                |                       |
       chi^2_min < 7.81             chi^2_min >= 7.81
       (95% CONFIRMED)                     |
       xi_best within 10%        +---------+----------+
       of 13/3 => UQFF locked    |                    |
                          7.81 <= chi^2 < 14.16   chi^2 >= 14.16
                          (TENSION, refit)        (FALSIFIED at 3-sigma)
\end{verbatim}

\section{5. CP4 Implementation}\label{cp4-implementation}

\begin{Shaded}
\begin{Highlighting}[]
\ImportTok{from}\NormalTok{ CondensedPhysics4 }\ImportTok{import}\NormalTok{ UQFF2027QuadrupleFalsifierCalculator}
\NormalTok{r }\OperatorTok{=}\NormalTok{ UQFF2027QuadrupleFalsifierCalculator().compute()}
\CommentTok{\# returns: xi\_best\_fit, chi2\_min, chi2\_threshold\_95cl\_3dof = 7.81,}
\CommentTok{\#          chi2\_threshold\_3sigma\_3dof = 14.16, status}
\end{Highlighting}
\end{Shaded}

For canonical inputs (all four observables at UQFF central values), the
calculator returns \(\chi^2_{\min} \approx 0\),
\(\xi_{\mathrm{best}} \approx 13/3\),
\texttt{status\ =\ "CONFIRMED\ (xi\ within\ 10\%\ of\ canonical\ 13/3)"}.

\section{6. Mutual Self-Locking Proof}\label{mutual-self-locking-proof}

A given \(\xi\) is feasible iff all four residuals are simultaneously
\(< 1\sigma\). Solving each constraint individually:

\[
\xi_{P6} = \xi_0, \qquad \xi_{P10} = \xi_0 \cdot \left(\frac{E_{cut}}{6.3}\right)^2,
\] \[
\xi_{P11} = \xi_0 \cdot \left(\frac{R_{21/22}}{0.144}\right)^{4}, \qquad
\xi_{P12} = \frac{\xi_0}{0.01525} \cdot \left(\frac{\sigma_8}{0.7851} - 1\right).
\]

Consistency requires \(|\xi_{P6} - \xi_{P10}| < 0.43\),
\(|\xi_{P6} - \xi_{P11}| < 0.43\), \(|\xi_{P6} - \xi_{P12}| < 0.43\) at
\(1\sigma\). The probability of accidental coincidence under random
alternative theories is \(\le 0.05^4 = 6.25 \times 10^{-6}\) at
\(2\sigma\).

\section{7. References}\label{references}

{[}1{]} PAPER\_1163 --- IceCube-Gen2 P10 neutrino Cherenkov cutoff.
{[}2{]} PAPER\_1174 --- UQFF closed ledger P6 KK tower forecast. {[}3{]}
PAPER\_1175 --- UQFF P11 LIGO O5 ringdown spectral offset. {[}4{]}
PAPER\_1176 --- UQFF P12 Euclid sigma\_8 R26 saturation. {[}5{]}
PAPER\_1177 --- UQFF 2027 joint falsifier triple (P6+P11+P12). {[}6{]}
PAPER\_1178 --- UQFF P13 DESI Y5 strict-static dark energy.


% Session 258 auto-injected bibliography stub
\bibliographystyle{unsrt}
\bibliography{refs}
\end{document}
